\documentclass[a4paper,12pt]{article}
\usepackage[brazil]{babel}
\usepackage[utf8]{inputenc}
\usepackage[T1]{fontenc}
\usepackage{amsmath}
\usepackage{array}
\usepackage{geometry}
\geometry{margin=2.5cm}

\begin{document}

\section*{Questão 121 — ENEM 2024 — Ciências da Natureza}

\subsection*{Interpretação e Estratégia de Resolução}

A questão solicita a compreensão da definição atual e científica do termo "compostos orgânicos", diferenciando-a do senso comum e das associações populares.

\textbf{Termos de comando:}
\begin{itemize}
  \item \textbf{Defina}: estabelecer significado correto com base em conceitos científicos.
  \item \textbf{Analise}: examinar as alternativas para identificar a correta.
\end{itemize}

\textbf{O que a questão pede?}
Reconhecer como a Química define compostos orgânicos, especialmente quanto à sua composição química.

\textbf{Estratégia para resposta:}
Distinguir definição tradicional, associada a seres vivos, da definição química, centrada no carbono. Analisar alternativas para identificar a que define composados orgânicos pela característica química fundamental, descartando confusões com origem e propriedades.


\subsection*{Conteúdo e Mini Revisão Teórica}

\begin{tabular}{|>{\raggedright}p{4.5cm}|p{7cm}|p{5cm}|}
\hline
\textbf{Conceito} & \textbf{Ideia Chave} & \textbf{Aplicação} \\
\hline
Compostos Orgânicos & Possuem carbono como elemento principal & Definição científica correta para classificar compostos orgânicos \\
\hline
Definição Tradicional x Científica & Uso popular associa a seres vivos; Química associa-se à composição química & Diferenciar conceito para correta interpretação \\
\hline
Propriedades e Origem & Não definem compostos orgânicos (ex: biodegradabilidade, origem, produção sem agrotóxicos) & Descartar alternativas que confundem propriedades com definição \\
\hline
\end{tabular}

\vspace{0.5cm}

\textbf{Teoria}\newline
Compostos orgânicos são definidos pela presença predominante de carbono em suas moléculas - base da Química Orgânica. Historicamente, acreditava-se que somente seres vivos podiam produzir esses compostos, mas hoje sabemos que podem ser sintetizados artificialmente. Assim, o termo engloba uma vasta gama de substâncias fundamentadas na composição química.


\subsection*{Resolução e \textit{Pulo do Gato}}

Inicialmente, o senso comum relaciona compostos orgânicos a seres vivos, mas essa noção é incompleta para fins científicos.

O carbono possui a capacidade única de formar longas cadeias e estruturas diversas, característica fundamental para classificar compostos orgânicos.

Portanto, compostos orgânicos são substâncias que contêm carbono como elemento principal.

Das alternativas, apenas a E apresenta essa característica corretamente; as demais vinculam o conceito a propriedades ou origens não definidoras.

\textbf{Pulo do Gato:} lembre-se sempre de que a classificação dos compostos orgânicos é definida pela presença de carbono em sua estrutura, não por sua origem, função ou processo produtivo.


\subsection*{Armadilhas e Distratores}

\begin{itemize}
  \item \textbf{A}: associar compostos orgânicos a benefícios à saúde; incorreto, pois definição é structural.
  \item \textbf{B}: confundir biodegradabilidade com definição química; nem todos são biodegradáveis.
  \item \textbf{C}: ligar compostos orgânicos só à síntese a partir de gás carbônico; definição independe da origem.
  \item \textbf{D}: confundir definição química com práticas agrícolas (produção sem agrotóxicos).
\end{itemize}


\subsection*{Código da Questão}

CN\_QUIMICA\_DEF\_001

\end{document}